% Adapted from an original template by Hlyni Arnórssyni, Reykjavik University, Iceland

\documentclass{article}

% Adapted from an original template by Hlyni Arnórssyni, Reykjavik University, Iceland
%
% ------------------------------ SETTINGS
\usepackage{geometry}

\geometry{
	paper=a4paper, % Paper size
	top=2.5cm, % Top margin
	bottom=2.5cm, % Bottom margin
	left=2.5cm, % Left margin
	right=2.4cm, % Right margin
	headheight=0.75cm, % Header height
	footskip=1.5cm, % Space from the bottom margin to the baseline of the footer
	headsep=0.75cm, % Space from the top margin to the baseline of the header
	%showframe, % Uncomment to show how the type block is set on the page
}

\usepackage{blindtext}
%-------------------------------- Character encoding ----------------------------
\usepackage[T1]{fontenc}
\usepackage[utf8]{inputenc}

%----------------------------- Mathematics packages from AMS ---------------

\usepackage{amsmath, amsfonts, amsthm, amssymb}
\usepackage{braket, nicefrac}

% ----------- International System of Units
\usepackage{siunitx}

%------------------------------ Lists / numbers -------------------------
\usepackage{enumitem, multicol}

%------------------------------- Figure insertions --------------
\usepackage{graphicx, float}  % Use option [H] to force the placement of a figure
\usepackage{keystroke}
\usepackage{pgfplots}\usepgfplotslibrary{units}\pgfplotsset{compat=1.16}

%------------------------------- Line Spacing --------------
\usepackage{setspace}

%------------------------------- Depth of the ToC --------------
\setcounter{tocdepth}{2}

%%%%%%%%%%%%%%%%%%%%%%%%%% Hyperlink References %%%%%%%%%%%%%%%%%%%%%%%%%%%
\usepackage{hyperref}

%--------------------% Storage Path for images %-----------------%
\graphicspath{{graphics/}{Graphics/}{./}}

%%%%%%%%%%%%%%%%%%%%%%%%%% Environments %%%%%%%%%%%%%%%%%%%%%%%%%%%
\renewenvironment{abstract}{
    \begin{center}
    \textbf{Abstract}
    \vspace{0.5cm}
    \par\itshape
    \begin{minipage}{0.8\linewidth}}{\end{minipage}
    \noindent\ignorespaces
    \end{center}
}

\newenvironment{keywords}{
    \begin{center}
    \textbf{Keywords}
    \vspace{0.5cm}
    \par
    \begin{minipage}{0.8\linewidth}}{\end{minipage}
    \noindent\ignorespaces
    \end{center}
}

\newenvironment{preface}{
    \begin{center}
    \textbf{Preface}
    \vspace{0.5cm}
    \par
    \begin{minipage}{0.8\linewidth}}{\end{minipage}
    \noindent\ignorespaces
    \end{center}
}

\newenvironment{acknowledgements}{
    \begin{center}
    \textbf{Acknowledgements}
    \vspace{0.5cm}
    \par
    \begin{minipage}{0.8\linewidth}}{\end{minipage}
    \noindent\ignorespaces
    \end{center}
}
\usepackage{listings}
\usepackage{lmodern}  % for bold teletype font
\usepackage{amsmath}  % for \hookrightarrow
\usepackage{xcolor}   % for \textcolor
\usepackage{caption}
\usepackage{subcaption}
\bibliographystyle{ieeetr}

\lstset{
  basicstyle=\ttfamily,
  columns=fullflexible,
  frame=single,
  breaklines=true
}

\begin{document}
%Title of the report, name of coworkers (of experiment and of report).
\begin{titlepage}
	\centering
	\includegraphics[width=0.6\textwidth]{Graphics/BI_logo.png}\par
	\vspace{5cm}

	{\scshape\huge Tardigrade's possible genes of interest\par} 
	\vspace{1cm}
	{\Large \today\par}
	\vfill
	
	%%%% PROJECT TITLE
	{\huge\bfseries Homework number 4\par}
	\vfill
	
	%%%% AUTHOR(S)
	{\Large\itshape by Sapozhnikov N.A.}\par
	\vspace{1.5cm}

	\vfill


	\vfill
\end{titlepage}

\newpage

\section{Abstract}
Tardigrades are the first known animal to survive in space. Thus we are trying to predict genes, responsible for utterly good DNA reparation, processing results of a TMS (Tandem Mass Spectrometry), looking for proteins that are most likely in a nucleus.

\section{Introduction}
Tardigrades, fascinatingly resilient microscopic animals also known as water bears or moss piglets, exhibit an unparalleled ability to endure extreme environmental conditions, ranging from the high altitudes of the Himalayas to the abyssal depths of the deep sea. Classified as "extremophiles," they can endure freezing temperatures, complete dehydration, high pressures exceeding 1,200 atmospheres, and radiation levels far surpassing those tolerated by other animals. Notably, tardigrades became the first known animals to survive the harsh conditions of space during the "Tardigrades In Space" (TARDIS) project in September 2007 \cite{Erdmann2017-bc}.

This research not only expands our understanding of extremophiles but also holds potential implications for fields such as radiation biology and genetic engineering. Unlocking the secrets encoded in the tardigrade genome may pave the way for innovative strategies to enhance stress resistance in other organisms and even shed light on mechanisms relevant to human health and longevity.

\section{Methods}
To process our TMS data we create a local database consisting from our short peptides. We used local alignment-based search with diamond \cite{buchfink_fast_2015} utility and blast with Protein-Protein BLAST 2.5.0. Then utilities TargetP \cite{armenteros_detecting_2019} and WoLF PSORT \cite{noauthor_wolf_nodate} were used to narrow down the amount of genes which are only take their effect in a nucleus. 

\section{Results}
Diamond blasting showed only 2 genes: g4106.t1 and g12510.t1. Niether of these genes were predicted to be in a nucleus by both TargetP and WoLF PSORT. Blastp found 118 genes 34 of which were unique. The analysis with TargetP was performed, list of genes marked as signal peptides (SP) was obtained. Genes from this list were also analized with WoLF PSORT.

\section{Discussion}
Discussion text

\newpage
\section{Bibliography}

\bibliography{references}
%\newpage
%\appendix
%\section{Appendix}
%Put data files, CAD drawings, additional sketches, etc.

\end{document}

